%% kapa: exam=1 round=19 year=2008 period=1 form=A subject=civil_law subject_ko=민법 seq=14 number=14
%% source: Q-Net 감정평가사 기출문제 (cst00309) artlSeq=1001962 file=extracted/1008/1001962/2008년 제19회 감정평가사 1차시험 기출문제/1교시.hwp
%% conversion: kordoc 4.24.51 -> tools/gen1.py (auto); review=proofread-sonnet; flags=none
\begin{question}{14}
취소와 무효에 관한 설명으로 옳은 것은? (다툼이 있는 경우에는 판례에 의함)
\begin{choices}
  \item 법률행위가 성립한 후 채무자의 귀책사유로 이행이 불가능하게 된 경우, 그 법률행위는 무효이다.
  \item 사회질서에 반하여 무효가 된 법률행위라도 이를 알고 추인하면 새로운 법률행위가 된다.
  \item 강박에 의한 의사표시의 취소는 악의의 제3자에 대하여도 대항할 수 없다.
  \item 토지소유자가 사기를 당하여 지상권을 설정한 후에 그 토지가 포괄승계된 경우, 포괄승계인은 지상권설정계약을 취소할 수 있다.
  \item 유동적 무효의 상태에 있는 거래계약의 당사자는 상대방의 협력의무 불이행을 이유로 일방적으로 유동적 무효의 상태에 있는 거래계약 자체를 해제할 수 있다.
\end{choices}
\end{question}
\answer{4}
